% ============================================================
%  macros.tex  --  Farben, Listings-Stile, tcolorbox-Umgebungen
%                  und Hilfsmakros fuer die TilVista-Dokumentation
%
%  Wird per % ============================================================
%  macros.tex  --  Farben, Listings-Stile, tcolorbox-Umgebungen
%                  und Hilfsmakros fuer die TilVista-Dokumentation
%
%  Wird per % ============================================================
%  macros.tex  --  Farben, Listings-Stile, tcolorbox-Umgebungen
%                  und Hilfsmakros fuer die TilVista-Dokumentation
%
%  Wird per % ============================================================
%  macros.tex  --  Farben, Listings-Stile, tcolorbox-Umgebungen
%                  und Hilfsmakros fuer die TilVista-Dokumentation
%
%  Wird per \input{src/packages/macros} in tilvista_doku.tex
%  eingebunden, nach packages.tex und vor \begin{document}.
% ============================================================

% ============================================================
%  FARBEN
% ============================================================

% Icon Farbschema - TilVista
\definecolor{tvwhite}	{HTML}{d6d6d6}
\definecolor{tvpink1}	{HTML}{d8bfcd}
\definecolor{tvpink2}	{HTML}{daa9c8}
\definecolor{tvpink3}	{HTML}{dc8bc0}
\definecolor{tvpurple1}	{HTML}{a3add6}
\definecolor{tvpurple2}	{HTML}{646ec4}

% Hauptfarbe fuer C++-Schluesselwoerter und Links
\definecolor{cppblue}  {HTML}{003E7E}

% Qt-typisches Gruen (Signale, Qt-spezifische Klassen)
\definecolor{qtgreen}  {HTML}{2E7D32}

% Hintergrund fuer Code-Bloecke (sehr helles Grau-Blau)
\definecolor{codebg}   {HTML}{F4F6F9}

% Rahmenfarbe fuer Code-Bloecke
\definecolor{framerule}{HTML}{B0BEC5}

% Hinweis-Box (gelb)
\definecolor{calloutbg}{HTML}{FFF8E1}
\definecolor{calloutfr}{HTML}{F9A825}

% Warnungs-Box (rot)
\definecolor{warnbg}   {HTML}{FFEBEE}
\definecolor{warnfr}   {HTML}{C62828}

% Info-Box / Qt-Hintergrund (blau)
\definecolor{infobg}   {HTML}{E3F2FD}
\definecolor{infofr}   {HTML}{1565C0}

% Neue-Funktion-Box (gruen)
\definecolor{newbg}    {HTML}{E8F5E9}
\definecolor{newfr}    {HTML}{2E7D32}

% Link-Farbe (dunkles Marineblau)
\definecolor{linkblue} {HTML}{1A237E}

% CMake-Schluesselwoerter (Dunkelrot)
\definecolor{cmakered} {HTML}{8B1A1A}

% ============================================================
%  LISTINGS-STILE
% ============================================================

% -- Basis-Stil fuer C++ -------------------------------------
\lstdefinestyle{cpp}{
  language        = C++,
  basicstyle      = \ttfamily\small,
  % Schluesselwoerter fett + blau
  keywordstyle    = \bfseries\color{cppblue},
  % Kommentare kursiv + hellgrau
  commentstyle    = \itshape\color{gray!65},
  % Strings dunkelrot
  stringstyle     = \color{BrickRed},
  % Zeilennummern (klein, grau, links)
  numberstyle     = \tiny\color{gray!60},
  numbers         = left,
  numbersep       = 9pt,
  % Hintergrund und Rahmen
  backgroundcolor = \color{codebg},
  frame           = single,
  framerule       = 0.5pt,
  rulecolor       = \color{framerule},
  % Zeilenumbruch bei langen Zeilen
  breaklines      = true,
  breakatwhitespace = true,
  tabsize         = 4,
  showstringspaces = false,
  % WICHTIG: extendedchars=false verhindert, dass listings
  % versucht, UTF-8-Bytes selbst zu interpretieren.
  % Umlaute in lstlisting-Bloecken muessen als ASCII-Umschreibung
  % stehen (ae, oe, ue, ss). Im normalen Fliestext ausserhalb
  % von lstlisting funktionieren Umlaute dank inputenc normal.
  extendedchars   = false,
  % Zusaetzliche Qt- und C++-Schluesselwoerter
  morekeywords    = {signals,slots,emit,Q_OBJECT,override,
                     nullptr,auto,constexpr,explicit,
                     QString,QWidget,QTimer,QPixmap,QLabel,
                     QMainWindow,QObject,QThread,Qt,
                     bool,int,double,float,void},
  xleftmargin     = 2em,
  xrightmargin    = 0.5em,
  aboveskip       = 6pt,
  belowskip       = 6pt
}

% -- Stil fuer Header-Dateien (.h) ---------------------------
% Identisch mit cpp, aber mit leicht blaulichem Hintergrund,
% um Header optisch von .cpp-Dateien zu unterscheiden.
\lstdefinestyle{header}{
  style           = cpp,
  backgroundcolor = \color{blue!5}
}

% -- Stil fuer CMake-Dateien ---------------------------------
\lstdefinestyle{cmake}{
  language        = {},       % kein eingebautes CMake-Highlighting
  basicstyle      = \ttfamily\small,
  keywordstyle    = \bfseries\color{cmakered},
  commentstyle    = \itshape\color{gray!65},
  stringstyle     = \color{BrickRed},
  numberstyle     = \tiny\color{gray!60},
  numbers         = left,
  numbersep       = 9pt,
  backgroundcolor = \color{codebg},
  frame           = single,
  framerule       = 0.5pt,
  rulecolor       = \color{framerule},
  breaklines      = true,
  tabsize         = 2,
  showstringspaces = false,
  extendedchars   = false,
  morekeywords    = {cmake_minimum_required,project,set,
                     find_package,add_executable,
                     target_link_libraries,target_compile_definitions,
                     target_include_directories,
                     set_target_properties,
                     add_custom_command,add_compile_definitions,
                     function,endfunction,if,else,endif,
                     find_program,message,include,install},
  xleftmargin     = 2em,
  xrightmargin    = 0.5em,
  aboveskip       = 6pt,
  belowskip       = 6pt
}

% Standard-Stil fuer alle \begin{lstlisting} ohne style=...
\lstset{style=cpp}

% ============================================================
%  tcolorbox-UMGEBUNGEN (farbige Infokaesten)
% ============================================================

% -- Hinweis (gelb) ------------------------------------------
% Fuer allgemeine Tipps und Erklaerungen
\newtcolorbox{hinweis}[1][]{
  enhanced, breakable,
  colback    = calloutbg,
  colframe   = calloutfr,
  fonttitle  = \bfseries\small,
  title      = {Hinweis},
  left=6pt, right=6pt, top=4pt, bottom=4pt,
  #1
}

% -- Warnung (rot) -------------------------------------------
% Fuer Fallstricke, undefiniertes Verhalten, Deadlock-Gefahren
\newtcolorbox{warnung}[1][]{
  enhanced, breakable,
  colback    = warnbg,
  colframe   = warnfr,
  fonttitle  = \bfseries\small,
  title      = {Achtung},
  left=6pt, right=6pt, top=4pt, bottom=4pt,
  #1
}

% -- Qt-Hintergrund (blau) -----------------------------------
% Fuer Qt-spezifische Konzepterklarungen mit Doku-Links
\newtcolorbox{qtinfo}[1][]{
  enhanced, breakable,
  colback    = infobg,
  colframe   = infofr,
  fonttitle  = \bfseries\small,
  title      = {Qt-Hintergrund},
  left=6pt, right=6pt, top=4pt, bottom=4pt,
  #1
}

% -- Neu in v0.5.42 (gruen) ----------------------------------
% Hebt Aenderungen der aktuellen Version hervor
\newtcolorbox{neuin}[1][]{
  enhanced, breakable,
  colback    = newbg,
  colframe   = newfr,
  fonttitle  = \bfseries\small,
  title      = {Neu in v0.5.42},
  left=6pt, right=6pt, top=4pt, bottom=4pt,
  #1
}

% ============================================================
%  HILFSMAKROS (Semantische Auszeichnung)
% ============================================================

% Dateinamen: monospace, z.B. \datei{mainwindow.h}
\newcommand{\datei}[1]{\texttt{#1}}

% Klassennamen: monospace + fett, z.B. \klasse{MainWindow}
\newcommand{\klasse}[1]{\texttt{\textbf{#1}}}

% Methoden mit Klammern: \methode{buildUi} --> buildUi()
\newcommand{\methode}[1]{\texttt{#1()}}

% Member-Variablen: \member{m\_dirBar}
\newcommand{\member}[1]{\texttt{#1}}

% Signale (Qt-gruen): \signal{dirLoaded}
\newcommand{\signal}[1]{\texttt{\color{qtgreen}#1}}

% Slots (blau): \slot{onDirChanged}
\newcommand{\slot}[1]{\texttt{\color{cppblue}#1}}

% Praeprozessor-Makros (dunkelrot): \makro{TV\_NO\_MULTIMEDIA}
\newcommand{\makro}[1]{\texttt{\color{BrickRed}#1}}

% CMake-Befehle: \cmake{find\_package}
\newcommand{\cmake}[1]{\texttt{\color{cmakered}#1}}

% Namespace-Praefix: \ns{TV} --> TV::
\newcommand{\ns}[1]{\texttt{#1::}}

% Versionsnummer hervorgehoben: \version{0.5.42}
\newcommand{\version}[1]{\textbf{v#1}}

% ============================================================
%  SEITENSTIL (fancyhdr)
% ============================================================

\pagestyle{fancy}
\fancyhf{}  % alle Felder leeren

% Linke Seite (gerade Seitenzahl): Icon + Kapitelname
\fancyhead[LE]{%
  \raisebox{-2pt}{\includegraphics[height=11pt]{graphics/icon}}%
  \quad\small\leftmark%
}
% Rechte Seite (ungerade Seitenzahl): Abschnittsname
\fancyhead[RO]{\small\rightmark}
% Fusszeile zentriert: Seitenzahl
\fancyfoot[C]{\small\thepage}

\renewcommand{\headrulewidth}{0.4pt}
\renewcommand{\footrulewidth}{0pt}
 in tilvista_doku.tex
%  eingebunden, nach packages.tex und vor \begin{document}.
% ============================================================

% ============================================================
%  FARBEN
% ============================================================

% Icon Farbschema - TilVista
\definecolor{tvwhite}	{HTML}{d6d6d6}
\definecolor{tvpink1}	{HTML}{d8bfcd}
\definecolor{tvpink2}	{HTML}{daa9c8}
\definecolor{tvpink3}	{HTML}{dc8bc0}
\definecolor{tvpurple1}	{HTML}{a3add6}
\definecolor{tvpurple2}	{HTML}{646ec4}

% Hauptfarbe fuer C++-Schluesselwoerter und Links
\definecolor{cppblue}  {HTML}{003E7E}

% Qt-typisches Gruen (Signale, Qt-spezifische Klassen)
\definecolor{qtgreen}  {HTML}{2E7D32}

% Hintergrund fuer Code-Bloecke (sehr helles Grau-Blau)
\definecolor{codebg}   {HTML}{F4F6F9}

% Rahmenfarbe fuer Code-Bloecke
\definecolor{framerule}{HTML}{B0BEC5}

% Hinweis-Box (gelb)
\definecolor{calloutbg}{HTML}{FFF8E1}
\definecolor{calloutfr}{HTML}{F9A825}

% Warnungs-Box (rot)
\definecolor{warnbg}   {HTML}{FFEBEE}
\definecolor{warnfr}   {HTML}{C62828}

% Info-Box / Qt-Hintergrund (blau)
\definecolor{infobg}   {HTML}{E3F2FD}
\definecolor{infofr}   {HTML}{1565C0}

% Neue-Funktion-Box (gruen)
\definecolor{newbg}    {HTML}{E8F5E9}
\definecolor{newfr}    {HTML}{2E7D32}

% Link-Farbe (dunkles Marineblau)
\definecolor{linkblue} {HTML}{1A237E}

% CMake-Schluesselwoerter (Dunkelrot)
\definecolor{cmakered} {HTML}{8B1A1A}

% ============================================================
%  LISTINGS-STILE
% ============================================================

% -- Basis-Stil fuer C++ -------------------------------------
\lstdefinestyle{cpp}{
  language        = C++,
  basicstyle      = \ttfamily\small,
  % Schluesselwoerter fett + blau
  keywordstyle    = \bfseries\color{cppblue},
  % Kommentare kursiv + hellgrau
  commentstyle    = \itshape\color{gray!65},
  % Strings dunkelrot
  stringstyle     = \color{BrickRed},
  % Zeilennummern (klein, grau, links)
  numberstyle     = \tiny\color{gray!60},
  numbers         = left,
  numbersep       = 9pt,
  % Hintergrund und Rahmen
  backgroundcolor = \color{codebg},
  frame           = single,
  framerule       = 0.5pt,
  rulecolor       = \color{framerule},
  % Zeilenumbruch bei langen Zeilen
  breaklines      = true,
  breakatwhitespace = true,
  tabsize         = 4,
  showstringspaces = false,
  % WICHTIG: extendedchars=false verhindert, dass listings
  % versucht, UTF-8-Bytes selbst zu interpretieren.
  % Umlaute in lstlisting-Bloecken muessen als ASCII-Umschreibung
  % stehen (ae, oe, ue, ss). Im normalen Fliestext ausserhalb
  % von lstlisting funktionieren Umlaute dank inputenc normal.
  extendedchars   = false,
  % Zusaetzliche Qt- und C++-Schluesselwoerter
  morekeywords    = {signals,slots,emit,Q_OBJECT,override,
                     nullptr,auto,constexpr,explicit,
                     QString,QWidget,QTimer,QPixmap,QLabel,
                     QMainWindow,QObject,QThread,Qt,
                     bool,int,double,float,void},
  xleftmargin     = 2em,
  xrightmargin    = 0.5em,
  aboveskip       = 6pt,
  belowskip       = 6pt
}

% -- Stil fuer Header-Dateien (.h) ---------------------------
% Identisch mit cpp, aber mit leicht blaulichem Hintergrund,
% um Header optisch von .cpp-Dateien zu unterscheiden.
\lstdefinestyle{header}{
  style           = cpp,
  backgroundcolor = \color{blue!5}
}

% -- Stil fuer CMake-Dateien ---------------------------------
\lstdefinestyle{cmake}{
  language        = {},       % kein eingebautes CMake-Highlighting
  basicstyle      = \ttfamily\small,
  keywordstyle    = \bfseries\color{cmakered},
  commentstyle    = \itshape\color{gray!65},
  stringstyle     = \color{BrickRed},
  numberstyle     = \tiny\color{gray!60},
  numbers         = left,
  numbersep       = 9pt,
  backgroundcolor = \color{codebg},
  frame           = single,
  framerule       = 0.5pt,
  rulecolor       = \color{framerule},
  breaklines      = true,
  tabsize         = 2,
  showstringspaces = false,
  extendedchars   = false,
  morekeywords    = {cmake_minimum_required,project,set,
                     find_package,add_executable,
                     target_link_libraries,target_compile_definitions,
                     target_include_directories,
                     set_target_properties,
                     add_custom_command,add_compile_definitions,
                     function,endfunction,if,else,endif,
                     find_program,message,include,install},
  xleftmargin     = 2em,
  xrightmargin    = 0.5em,
  aboveskip       = 6pt,
  belowskip       = 6pt
}

% Standard-Stil fuer alle \begin{lstlisting} ohne style=...
\lstset{style=cpp}

% ============================================================
%  tcolorbox-UMGEBUNGEN (farbige Infokaesten)
% ============================================================

% -- Hinweis (gelb) ------------------------------------------
% Fuer allgemeine Tipps und Erklaerungen
\newtcolorbox{hinweis}[1][]{
  enhanced, breakable,
  colback    = calloutbg,
  colframe   = calloutfr,
  fonttitle  = \bfseries\small,
  title      = {Hinweis},
  left=6pt, right=6pt, top=4pt, bottom=4pt,
  #1
}

% -- Warnung (rot) -------------------------------------------
% Fuer Fallstricke, undefiniertes Verhalten, Deadlock-Gefahren
\newtcolorbox{warnung}[1][]{
  enhanced, breakable,
  colback    = warnbg,
  colframe   = warnfr,
  fonttitle  = \bfseries\small,
  title      = {Achtung},
  left=6pt, right=6pt, top=4pt, bottom=4pt,
  #1
}

% -- Qt-Hintergrund (blau) -----------------------------------
% Fuer Qt-spezifische Konzepterklarungen mit Doku-Links
\newtcolorbox{qtinfo}[1][]{
  enhanced, breakable,
  colback    = infobg,
  colframe   = infofr,
  fonttitle  = \bfseries\small,
  title      = {Qt-Hintergrund},
  left=6pt, right=6pt, top=4pt, bottom=4pt,
  #1
}

% -- Neu in v0.5.42 (gruen) ----------------------------------
% Hebt Aenderungen der aktuellen Version hervor
\newtcolorbox{neuin}[1][]{
  enhanced, breakable,
  colback    = newbg,
  colframe   = newfr,
  fonttitle  = \bfseries\small,
  title      = {Neu in v0.5.42},
  left=6pt, right=6pt, top=4pt, bottom=4pt,
  #1
}

% ============================================================
%  HILFSMAKROS (Semantische Auszeichnung)
% ============================================================

% Dateinamen: monospace, z.B. \datei{mainwindow.h}
\newcommand{\datei}[1]{\texttt{#1}}

% Klassennamen: monospace + fett, z.B. \klasse{MainWindow}
\newcommand{\klasse}[1]{\texttt{\textbf{#1}}}

% Methoden mit Klammern: \methode{buildUi} --> buildUi()
\newcommand{\methode}[1]{\texttt{#1()}}

% Member-Variablen: \member{m\_dirBar}
\newcommand{\member}[1]{\texttt{#1}}

% Signale (Qt-gruen): \signal{dirLoaded}
\newcommand{\signal}[1]{\texttt{\color{qtgreen}#1}}

% Slots (blau): \slot{onDirChanged}
\newcommand{\slot}[1]{\texttt{\color{cppblue}#1}}

% Praeprozessor-Makros (dunkelrot): \makro{TV\_NO\_MULTIMEDIA}
\newcommand{\makro}[1]{\texttt{\color{BrickRed}#1}}

% CMake-Befehle: \cmake{find\_package}
\newcommand{\cmake}[1]{\texttt{\color{cmakered}#1}}

% Namespace-Praefix: \ns{TV} --> TV::
\newcommand{\ns}[1]{\texttt{#1::}}

% Versionsnummer hervorgehoben: \version{0.5.42}
\newcommand{\version}[1]{\textbf{v#1}}

% ============================================================
%  SEITENSTIL (fancyhdr)
% ============================================================

\pagestyle{fancy}
\fancyhf{}  % alle Felder leeren

% Linke Seite (gerade Seitenzahl): Icon + Kapitelname
\fancyhead[LE]{%
  \raisebox{-2pt}{\includegraphics[height=11pt]{graphics/icon}}%
  \quad\small\leftmark%
}
% Rechte Seite (ungerade Seitenzahl): Abschnittsname
\fancyhead[RO]{\small\rightmark}
% Fusszeile zentriert: Seitenzahl
\fancyfoot[C]{\small\thepage}

\renewcommand{\headrulewidth}{0.4pt}
\renewcommand{\footrulewidth}{0pt}
 in tilvista_doku.tex
%  eingebunden, nach packages.tex und vor \begin{document}.
% ============================================================

% ============================================================
%  FARBEN
% ============================================================

% Icon Farbschema - TilVista
\definecolor{tvwhite}	{HTML}{d6d6d6}
\definecolor{tvpink1}	{HTML}{d8bfcd}
\definecolor{tvpink2}	{HTML}{daa9c8}
\definecolor{tvpink3}	{HTML}{dc8bc0}
\definecolor{tvpurple1}	{HTML}{a3add6}
\definecolor{tvpurple2}	{HTML}{646ec4}

% Hauptfarbe fuer C++-Schluesselwoerter und Links
\definecolor{cppblue}  {HTML}{003E7E}

% Qt-typisches Gruen (Signale, Qt-spezifische Klassen)
\definecolor{qtgreen}  {HTML}{2E7D32}

% Hintergrund fuer Code-Bloecke (sehr helles Grau-Blau)
\definecolor{codebg}   {HTML}{F4F6F9}

% Rahmenfarbe fuer Code-Bloecke
\definecolor{framerule}{HTML}{B0BEC5}

% Hinweis-Box (gelb)
\definecolor{calloutbg}{HTML}{FFF8E1}
\definecolor{calloutfr}{HTML}{F9A825}

% Warnungs-Box (rot)
\definecolor{warnbg}   {HTML}{FFEBEE}
\definecolor{warnfr}   {HTML}{C62828}

% Info-Box / Qt-Hintergrund (blau)
\definecolor{infobg}   {HTML}{E3F2FD}
\definecolor{infofr}   {HTML}{1565C0}

% Neue-Funktion-Box (gruen)
\definecolor{newbg}    {HTML}{E8F5E9}
\definecolor{newfr}    {HTML}{2E7D32}

% Link-Farbe (dunkles Marineblau)
\definecolor{linkblue} {HTML}{1A237E}

% CMake-Schluesselwoerter (Dunkelrot)
\definecolor{cmakered} {HTML}{8B1A1A}

% ============================================================
%  LISTINGS-STILE
% ============================================================

% -- Basis-Stil fuer C++ -------------------------------------
\lstdefinestyle{cpp}{
  language        = C++,
  basicstyle      = \ttfamily\small,
  % Schluesselwoerter fett + blau
  keywordstyle    = \bfseries\color{cppblue},
  % Kommentare kursiv + hellgrau
  commentstyle    = \itshape\color{gray!65},
  % Strings dunkelrot
  stringstyle     = \color{BrickRed},
  % Zeilennummern (klein, grau, links)
  numberstyle     = \tiny\color{gray!60},
  numbers         = left,
  numbersep       = 9pt,
  % Hintergrund und Rahmen
  backgroundcolor = \color{codebg},
  frame           = single,
  framerule       = 0.5pt,
  rulecolor       = \color{framerule},
  % Zeilenumbruch bei langen Zeilen
  breaklines      = true,
  breakatwhitespace = true,
  tabsize         = 4,
  showstringspaces = false,
  % WICHTIG: extendedchars=false verhindert, dass listings
  % versucht, UTF-8-Bytes selbst zu interpretieren.
  % Umlaute in lstlisting-Bloecken muessen als ASCII-Umschreibung
  % stehen (ae, oe, ue, ss). Im normalen Fliestext ausserhalb
  % von lstlisting funktionieren Umlaute dank inputenc normal.
  extendedchars   = false,
  % Zusaetzliche Qt- und C++-Schluesselwoerter
  morekeywords    = {signals,slots,emit,Q_OBJECT,override,
                     nullptr,auto,constexpr,explicit,
                     QString,QWidget,QTimer,QPixmap,QLabel,
                     QMainWindow,QObject,QThread,Qt,
                     bool,int,double,float,void},
  xleftmargin     = 2em,
  xrightmargin    = 0.5em,
  aboveskip       = 6pt,
  belowskip       = 6pt
}

% -- Stil fuer Header-Dateien (.h) ---------------------------
% Identisch mit cpp, aber mit leicht blaulichem Hintergrund,
% um Header optisch von .cpp-Dateien zu unterscheiden.
\lstdefinestyle{header}{
  style           = cpp,
  backgroundcolor = \color{blue!5}
}

% -- Stil fuer CMake-Dateien ---------------------------------
\lstdefinestyle{cmake}{
  language        = {},       % kein eingebautes CMake-Highlighting
  basicstyle      = \ttfamily\small,
  keywordstyle    = \bfseries\color{cmakered},
  commentstyle    = \itshape\color{gray!65},
  stringstyle     = \color{BrickRed},
  numberstyle     = \tiny\color{gray!60},
  numbers         = left,
  numbersep       = 9pt,
  backgroundcolor = \color{codebg},
  frame           = single,
  framerule       = 0.5pt,
  rulecolor       = \color{framerule},
  breaklines      = true,
  tabsize         = 2,
  showstringspaces = false,
  extendedchars   = false,
  morekeywords    = {cmake_minimum_required,project,set,
                     find_package,add_executable,
                     target_link_libraries,target_compile_definitions,
                     target_include_directories,
                     set_target_properties,
                     add_custom_command,add_compile_definitions,
                     function,endfunction,if,else,endif,
                     find_program,message,include,install},
  xleftmargin     = 2em,
  xrightmargin    = 0.5em,
  aboveskip       = 6pt,
  belowskip       = 6pt
}

% Standard-Stil fuer alle \begin{lstlisting} ohne style=...
\lstset{style=cpp}

% ============================================================
%  tcolorbox-UMGEBUNGEN (farbige Infokaesten)
% ============================================================

% -- Hinweis (gelb) ------------------------------------------
% Fuer allgemeine Tipps und Erklaerungen
\newtcolorbox{hinweis}[1][]{
  enhanced, breakable,
  colback    = calloutbg,
  colframe   = calloutfr,
  fonttitle  = \bfseries\small,
  title      = {Hinweis},
  left=6pt, right=6pt, top=4pt, bottom=4pt,
  #1
}

% -- Warnung (rot) -------------------------------------------
% Fuer Fallstricke, undefiniertes Verhalten, Deadlock-Gefahren
\newtcolorbox{warnung}[1][]{
  enhanced, breakable,
  colback    = warnbg,
  colframe   = warnfr,
  fonttitle  = \bfseries\small,
  title      = {Achtung},
  left=6pt, right=6pt, top=4pt, bottom=4pt,
  #1
}

% -- Qt-Hintergrund (blau) -----------------------------------
% Fuer Qt-spezifische Konzepterklarungen mit Doku-Links
\newtcolorbox{qtinfo}[1][]{
  enhanced, breakable,
  colback    = infobg,
  colframe   = infofr,
  fonttitle  = \bfseries\small,
  title      = {Qt-Hintergrund},
  left=6pt, right=6pt, top=4pt, bottom=4pt,
  #1
}

% -- Neu in v0.5.42 (gruen) ----------------------------------
% Hebt Aenderungen der aktuellen Version hervor
\newtcolorbox{neuin}[1][]{
  enhanced, breakable,
  colback    = newbg,
  colframe   = newfr,
  fonttitle  = \bfseries\small,
  title      = {Neu in v0.5.42},
  left=6pt, right=6pt, top=4pt, bottom=4pt,
  #1
}

% ============================================================
%  HILFSMAKROS (Semantische Auszeichnung)
% ============================================================

% Dateinamen: monospace, z.B. \datei{mainwindow.h}
\newcommand{\datei}[1]{\texttt{#1}}

% Klassennamen: monospace + fett, z.B. \klasse{MainWindow}
\newcommand{\klasse}[1]{\texttt{\textbf{#1}}}

% Methoden mit Klammern: \methode{buildUi} --> buildUi()
\newcommand{\methode}[1]{\texttt{#1()}}

% Member-Variablen: \member{m\_dirBar}
\newcommand{\member}[1]{\texttt{#1}}

% Signale (Qt-gruen): \signal{dirLoaded}
\newcommand{\signal}[1]{\texttt{\color{qtgreen}#1}}

% Slots (blau): \slot{onDirChanged}
\newcommand{\slot}[1]{\texttt{\color{cppblue}#1}}

% Praeprozessor-Makros (dunkelrot): \makro{TV\_NO\_MULTIMEDIA}
\newcommand{\makro}[1]{\texttt{\color{BrickRed}#1}}

% CMake-Befehle: \cmake{find\_package}
\newcommand{\cmake}[1]{\texttt{\color{cmakered}#1}}

% Namespace-Praefix: \ns{TV} --> TV::
\newcommand{\ns}[1]{\texttt{#1::}}

% Versionsnummer hervorgehoben: \version{0.5.42}
\newcommand{\version}[1]{\textbf{v#1}}

% ============================================================
%  SEITENSTIL (fancyhdr)
% ============================================================

\pagestyle{fancy}
\fancyhf{}  % alle Felder leeren

% Linke Seite (gerade Seitenzahl): Icon + Kapitelname
\fancyhead[LE]{%
  \raisebox{-2pt}{\includegraphics[height=11pt]{graphics/icon}}%
  \quad\small\leftmark%
}
% Rechte Seite (ungerade Seitenzahl): Abschnittsname
\fancyhead[RO]{\small\rightmark}
% Fusszeile zentriert: Seitenzahl
\fancyfoot[C]{\small\thepage}

\renewcommand{\headrulewidth}{0.4pt}
\renewcommand{\footrulewidth}{0pt}
 in tilvista_doku.tex
%  eingebunden, nach packages.tex und vor \begin{document}.
% ============================================================

% ============================================================
%  FARBEN
% ============================================================

% Icon Farbschema - TilVista
\definecolor{tvwhite}	{HTML}{d6d6d6}
\definecolor{tvpink1}	{HTML}{d8bfcd}
\definecolor{tvpink2}	{HTML}{daa9c8}
\definecolor{tvpink3}	{HTML}{dc8bc0}
\definecolor{tvpurple1}	{HTML}{a3add6}
\definecolor{tvpurple2}	{HTML}{646ec4}

% Hauptfarbe fuer C++-Schluesselwoerter und Links
\definecolor{cppblue}  {HTML}{003E7E}

% Qt-typisches Gruen (Signale, Qt-spezifische Klassen)
\definecolor{qtgreen}  {HTML}{2E7D32}

% Hintergrund fuer Code-Bloecke (sehr helles Grau-Blau)
\definecolor{codebg}   {HTML}{F4F6F9}

% Rahmenfarbe fuer Code-Bloecke
\definecolor{framerule}{HTML}{B0BEC5}

% Hinweis-Box (gelb)
\definecolor{calloutbg}{HTML}{FFF8E1}
\definecolor{calloutfr}{HTML}{F9A825}

% Warnungs-Box (rot)
\definecolor{warnbg}   {HTML}{FFEBEE}
\definecolor{warnfr}   {HTML}{C62828}

% Info-Box / Qt-Hintergrund (blau)
\definecolor{infobg}   {HTML}{E3F2FD}
\definecolor{infofr}   {HTML}{1565C0}

% Neue-Funktion-Box (gruen)
\definecolor{newbg}    {HTML}{E8F5E9}
\definecolor{newfr}    {HTML}{2E7D32}

% Link-Farbe (dunkles Marineblau)
\definecolor{linkblue} {HTML}{1A237E}

% CMake-Schluesselwoerter (Dunkelrot)
\definecolor{cmakered} {HTML}{8B1A1A}

% ============================================================
%  LISTINGS-STILE
% ============================================================

% -- Basis-Stil fuer C++ -------------------------------------
\lstdefinestyle{cpp}{
  language        = C++,
  basicstyle      = \ttfamily\small,
  % Schluesselwoerter fett + blau
  keywordstyle    = \bfseries\color{cppblue},
  % Kommentare kursiv + hellgrau
  commentstyle    = \itshape\color{gray!65},
  % Strings dunkelrot
  stringstyle     = \color{BrickRed},
  % Zeilennummern (klein, grau, links)
  numberstyle     = \tiny\color{gray!60},
  numbers         = left,
  numbersep       = 9pt,
  % Hintergrund und Rahmen
  backgroundcolor = \color{codebg},
  frame           = single,
  framerule       = 0.5pt,
  rulecolor       = \color{framerule},
  % Zeilenumbruch bei langen Zeilen
  breaklines      = true,
  breakatwhitespace = true,
  tabsize         = 4,
  showstringspaces = false,
  % WICHTIG: extendedchars=false verhindert, dass listings
  % versucht, UTF-8-Bytes selbst zu interpretieren.
  % Umlaute in lstlisting-Bloecken muessen als ASCII-Umschreibung
  % stehen (ae, oe, ue, ss). Im normalen Fliestext ausserhalb
  % von lstlisting funktionieren Umlaute dank inputenc normal.
  extendedchars   = false,
  % Zusaetzliche Qt- und C++-Schluesselwoerter
  morekeywords    = {signals,slots,emit,Q_OBJECT,override,
                     nullptr,auto,constexpr,explicit,
                     QString,QWidget,QTimer,QPixmap,QLabel,
                     QMainWindow,QObject,QThread,Qt,
                     bool,int,double,float,void},
  xleftmargin     = 2em,
  xrightmargin    = 0.5em,
  aboveskip       = 6pt,
  belowskip       = 6pt
}

% -- Stil fuer Header-Dateien (.h) ---------------------------
% Identisch mit cpp, aber mit leicht blaulichem Hintergrund,
% um Header optisch von .cpp-Dateien zu unterscheiden.
\lstdefinestyle{header}{
  style           = cpp,
  backgroundcolor = \color{blue!5}
}

% -- Stil fuer CMake-Dateien ---------------------------------
\lstdefinestyle{cmake}{
  language        = {},       % kein eingebautes CMake-Highlighting
  basicstyle      = \ttfamily\small,
  keywordstyle    = \bfseries\color{cmakered},
  commentstyle    = \itshape\color{gray!65},
  stringstyle     = \color{BrickRed},
  numberstyle     = \tiny\color{gray!60},
  numbers         = left,
  numbersep       = 9pt,
  backgroundcolor = \color{codebg},
  frame           = single,
  framerule       = 0.5pt,
  rulecolor       = \color{framerule},
  breaklines      = true,
  tabsize         = 2,
  showstringspaces = false,
  extendedchars   = false,
  morekeywords    = {cmake_minimum_required,project,set,
                     find_package,add_executable,
                     target_link_libraries,target_compile_definitions,
                     target_include_directories,
                     set_target_properties,
                     add_custom_command,add_compile_definitions,
                     function,endfunction,if,else,endif,
                     find_program,message,include,install},
  xleftmargin     = 2em,
  xrightmargin    = 0.5em,
  aboveskip       = 6pt,
  belowskip       = 6pt
}

% Standard-Stil fuer alle \begin{lstlisting} ohne style=...
\lstset{style=cpp}

% ============================================================
%  tcolorbox-UMGEBUNGEN (farbige Infokaesten)
% ============================================================

% -- Hinweis (gelb) ------------------------------------------
% Fuer allgemeine Tipps und Erklaerungen
\newtcolorbox{hinweis}[1][]{
  enhanced, breakable,
  colback    = calloutbg,
  colframe   = calloutfr,
  fonttitle  = \bfseries\small,
  title      = {Hinweis},
  left=6pt, right=6pt, top=4pt, bottom=4pt,
  #1
}

% -- Warnung (rot) -------------------------------------------
% Fuer Fallstricke, undefiniertes Verhalten, Deadlock-Gefahren
\newtcolorbox{warnung}[1][]{
  enhanced, breakable,
  colback    = warnbg,
  colframe   = warnfr,
  fonttitle  = \bfseries\small,
  title      = {Achtung},
  left=6pt, right=6pt, top=4pt, bottom=4pt,
  #1
}

% -- Qt-Hintergrund (blau) -----------------------------------
% Fuer Qt-spezifische Konzepterklarungen mit Doku-Links
\newtcolorbox{qtinfo}[1][]{
  enhanced, breakable,
  colback    = infobg,
  colframe   = infofr,
  fonttitle  = \bfseries\small,
  title      = {Qt-Hintergrund},
  left=6pt, right=6pt, top=4pt, bottom=4pt,
  #1
}

% -- Neu in v0.5.42 (gruen) ----------------------------------
% Hebt Aenderungen der aktuellen Version hervor
\newtcolorbox{neuin}[1][]{
  enhanced, breakable,
  colback    = newbg,
  colframe   = newfr,
  fonttitle  = \bfseries\small,
  title      = {Neu in v0.5.42},
  left=6pt, right=6pt, top=4pt, bottom=4pt,
  #1
}

% ============================================================
%  HILFSMAKROS (Semantische Auszeichnung)
% ============================================================

% Dateinamen: monospace, z.B. \datei{mainwindow.h}
\newcommand{\datei}[1]{\texttt{#1}}

% Klassennamen: monospace + fett, z.B. \klasse{MainWindow}
\newcommand{\klasse}[1]{\texttt{\textbf{#1}}}

% Methoden mit Klammern: \methode{buildUi} --> buildUi()
\newcommand{\methode}[1]{\texttt{#1()}}

% Member-Variablen: \member{m\_dirBar}
\newcommand{\member}[1]{\texttt{#1}}

% Signale (Qt-gruen): \signal{dirLoaded}
\newcommand{\signal}[1]{\texttt{\color{qtgreen}#1}}

% Slots (blau): \slot{onDirChanged}
\newcommand{\slot}[1]{\texttt{\color{cppblue}#1}}

% Praeprozessor-Makros (dunkelrot): \makro{TV\_NO\_MULTIMEDIA}
\newcommand{\makro}[1]{\texttt{\color{BrickRed}#1}}

% CMake-Befehle: \cmake{find\_package}
\newcommand{\cmake}[1]{\texttt{\color{cmakered}#1}}

% Namespace-Praefix: \ns{TV} --> TV::
\newcommand{\ns}[1]{\texttt{#1::}}

% Versionsnummer hervorgehoben: \version{0.5.42}
\newcommand{\version}[1]{\textbf{v#1}}

% ============================================================
%  SEITENSTIL (fancyhdr)
% ============================================================

\pagestyle{fancy}
\fancyhf{}  % alle Felder leeren

% Linke Seite (gerade Seitenzahl): Icon + Kapitelname
\fancyhead[LE]{%
  \raisebox{-2pt}{\includegraphics[height=11pt]{graphics/icon}}%
  \quad\small\leftmark%
}
% Rechte Seite (ungerade Seitenzahl): Abschnittsname
\fancyhead[RO]{\small\rightmark}
% Fusszeile zentriert: Seitenzahl
\fancyfoot[C]{\small\thepage}

\renewcommand{\headrulewidth}{0.4pt}
\renewcommand{\footrulewidth}{0pt}
